\documentclass[10pt,twocolumn]{article}

% ── Packages ──
\usepackage[utf8]{inputenc}
\usepackage[T1]{fontenc}
\usepackage{amsmath,amssymb,amsthm}
\usepackage{booktabs}
\usepackage{hyperref}
\usepackage{graphicx}
\usepackage{natbib}
\usepackage[margin=1in]{geometry}
\usepackage{enumitem}
\usepackage{listings}
\usepackage{makecell}
\usepackage{multirow}
\usepackage{balance}

% ── Theorem Environments ──
\newtheorem{theorem}{Theorem}
\newtheorem{lemma}[theorem]{Lemma}
\newtheorem{definition}[theorem]{Definition}

% ── Code Listings ──
\lstset{basicstyle=\ttfamily\footnotesize,breaklines=true,frame=single}

% ── Title ──
\title{Cross-Architecture KV Cache Quantization:\\Why No Single Method Wins}
\author{Muhamad Fajar Putranto\\TaxPrime / PrimeCore.id, Jakarta, Indonesia}
\date{}

\begin{document}
\maketitle

% ═══════════════════════════════════════════════════════════
\begin{abstract}
KV cache quantization is critical for deploying large language models on
resource-constrained devices. We present the first systematic
cross-architecture perplexity evaluation of five KV cache quantization
methods across three fundamentally different attention designs: MQA
(Gemma~4~E2B, 1~KV head), wide-GQA (Mistral~7B, 8~KV heads), and
narrow-GQA (Qwen2-7B, 4~KV heads).

Our key finding is that \textbf{no single method dominates across
architectures}. KIVI wins 6 of 9 cells in our evaluation matrix, but
catastrophically fails on Gemma MQA at 2-bit (PPL~470 vs FP16~28).
FajarQuant~v2---which combines calibrated PCA rotation with 15\% outlier
extraction---is the strongest non-KIVI method on MQA and narrow-GQA,
achieving 69.7\% lower PPL than TurboQuant on Qwen2-7B at 2-bit.
Architecture dependence is the dominant signal: the ``right'' method
depends on the KV head count, not the bit width.

We additionally contribute: (1)~a canonical R-$\alpha$.1 model-surgery
evaluation protocol enabling fair cross-method comparison;
(2)~\textbf{FajarQuant~v2}, an outlier-aware calibrated PCA method that
improves 63--81\% over v1 at 2-bit; and (3)~a native Fajar Lang
implementation with compile-time safety guarantees (\texttt{@device}
context, SE023 type rule) running 5$\times$ faster than the Python
reference.
\end{abstract}

% ═══════════════════════════════════════════════════════════
\section{Introduction}

\subsection{The KV Cache Memory Problem}

Transformer-based large language models require storing key-value (KV) caches for autoregressive generation. For a model with $L$ layers, $H$ heads, dimension $d$ per head, and sequence length $N$:
\begin{equation}
\text{Memory} = 2 \times L \times H \times N \times d \times \text{sizeof}(\text{dtype})
\end{equation}
For a 7B model ($L{=}32$, $H{=}32$, $d{=}128$) at $N{=}4096$ in FP16: \textbf{4~GB}. Embedded targets typically have 256~KB--8~MB SRAM. Aggressive quantization is essential.

\subsection{Prior Work}

\textbf{TurboQuant} \citep{zandieh2025turboquant} applies random orthogonal rotation followed by coordinate-wise Lloyd-Max quantization, achieving MSE within $2.7\times$ of the information-theoretic optimum. \textbf{KIVI} \citep{liu2024kivi} observes that keys have large magnitude variation across channels and values across tokens, using per-channel key and per-token value quantization.

\subsection{Contributions}

We identify three opportunities to improve upon prior work:
\begin{enumerate}[nosep]
\item \textbf{The rotation is data-agnostic.} KV cache vectors exhibit moderate low-rank structure (SV1/SV2 ratio 2.73 on Gemma~4). PCA rotation can exploit this.
\item \textbf{Dequantization wastes memory.} Standard attention requires fully dequantizing each key, allocating $O(N \cdot d)$ extra memory.
\item \textbf{Flat bit allocation wastes budget.} Recent tokens receive far more attention weight than old tokens.
\end{enumerate}

% ═══════════════════════════════════════════════════════════
\section{Background: TurboQuant}

\begin{lemma}[Beta Distribution after Rotation]
For $x$ on the unit sphere $\mathbb{S}^{d-1}$, after orthogonal rotation $y = \Pi x$, each coordinate $y_i$ follows $\text{Beta}(\frac{d-1}{2}, \frac{d-1}{2})$ on $[-1, 1]$.
\end{lemma}

\begin{theorem}[TurboQuant Distortion Bound]
\begin{equation}
D_{\text{mse}}(b) \leq \frac{\sqrt{3}\pi}{2} \cdot \frac{1}{4^b} \leq 2.72 \cdot D^*(b)
\end{equation}
\end{theorem}

% ═══════════════════════════════════════════════════════════
\section{Evaluation Methodology}
\label{sec:methodology}

\subsection{Canonical Perplexity Protocol (R-$\alpha$.1)}

Prior KV cache quantization papers use varying evaluation protocols,
making cross-method comparison unreliable. We adopt the
\textbf{R-$\alpha$.1 model surgery} protocol, following
KVQuant~\citep{hooper2024kvquant} and SKVQ~\citep{duanmu2024skvq}:

\begin{enumerate}[nosep]
\item Load the full-precision model.
\item Patch the attention forward pass to apply quantization
      \emph{after} key/value projection but \emph{before} the
      attention score computation. This is ``model surgery''---the
      KV cache is quantized in-flight, not post-hoc.
\item Evaluate perplexity on WikiText-2 test using non-overlapping
      chunks of length $L{=}2048$, computing cross-entropy on
      positions $[1, L)$ per chunk.
\item Report perplexity as $\exp(\text{mean NLL})$ over $n{=}30$
      chunks (61,410 tokens).
\end{enumerate}

The FP16 baseline is produced by the same forward pass with
quantization disabled (identity patch). This ensures the baseline
reflects the model's actual WikiText-2 perplexity, not an artifact
of the evaluation harness.

\subsection{The v1 $\to$ v2 Transition}
\label{sec:v1v2}

FajarQuant v1 used PCA rotation calibrated on key data, with symmetric
per-tensor quantization. Under the original non-canonical evaluation
protocol (prefix+target post-hoc cache mutation), v1 appeared to
outperform all baselines at 2--3~bit on Gemma~4~E2B. When we switched
to the canonical R-$\alpha$.1 protocol, v1's advantage disappeared:
\textbf{v1 lost to TurboQuant by 5.6$\times$} on Gemma 2-bit
(PPL~125 vs PPL~40).

Diagnosis revealed two root causes: (1)~v1's per-tensor symmetric
quantization destroyed outlier channels that carry disproportionate
information in real KV cache; (2)~the PCA rotation without outlier
handling amplified rather than reduced quantization error on high-kurtosis
distributions.

FajarQuant v2 addresses both: it adds \textbf{15\% outlier extraction}
(top channels preserved in higher precision) and uses a \textbf{fused
Hadamard+quantize kernel} that spreads outliers before quantization.
V2 improves over v1 by 63--81\% at 2-bit across all three models
(Table~\ref{tab:ppl_crossmodel}).

% ═══════════════════════════════════════════════════════════
\section{Innovation 1: Adaptive PCA Rotation}

\subsection{Motivation}

Random rotation uniformizes coordinate distributions---optimal for worst-case data but suboptimal for structured KV cache vectors. We analyze real KV cache from Google Gemma~4~E2B and find moderate low-rank structure: the mean singular value ratio $\sigma_1/\sigma_2 = 2.73 \pm 0.56$ across 150 layer-head samples, with Layer~1 showing the strongest structure ($\sigma_1/\sigma_2 = 4.04$). The top-3 singular values capture 66.7\% of total energy in keys---significantly more than the 50\% expected from random data.

\subsection{Method}

Given calibration data $\{x_1, \ldots, x_N\}$, we compute the covariance $C = \frac{1}{N} \sum_i x_i x_i^T$ and extract eigenvectors via eigendecomposition to form $R_{\text{pca}}$. After PCA rotation, each coordinate corresponds to a principal component with distinct variance. We then apply per-coordinate uniform quantization with independent scale and zero-point per dimension---allocating quantization resolution proportional to each component's dynamic range.

\begin{theorem}[Adaptive Rotation Bound]
Let $C$ have eigenvalues $\lambda_1 \geq \cdots \geq \lambda_d$. With per-coordinate quantization after PCA rotation:
\begin{equation}
\text{MSE}(R_{\text{pca}}) \leq \text{MSE}(R_{\text{rand}})
\end{equation}
\end{theorem}

\subsection{Results on Real Gemma~4 Data}

We evaluate on KV cache extracted from Google Gemma~4~E2B (35 layers, 1~KV head, $d{=}256$) using 50 diverse prompts from WikiText-2. All methods use per-coordinate quantization for fair comparison.

\begin{table}[t]
\centering
\small
\caption{Key quantization MSE on real Gemma~4 E2B KV cache ($d{=}256$, 150 samples)}
\label{tab:mse_real}
\begin{tabular}{lccc}
\toprule
Bits & FajarQuant & KIVI & TurboQuant \\
\midrule
$b{=}2$ & $\mathbf{8.72 \times 10^{-4}}$ & $9.31 \times 10^{-4}$ & $9.18 \times 10^{-4}$ \\
$b{=}3$ & $\mathbf{1.60 \times 10^{-4}}$ & $1.66 \times 10^{-4}$ & $1.67 \times 10^{-4}$ \\
$b{=}4$ & $\mathbf{3.46 \times 10^{-5}}$ & $3.60 \times 10^{-5}$ & $3.63 \times 10^{-5}$ \\
\midrule
\multicolumn{4}{l}{\textit{Improvement over baselines:}} \\
$b{=}2$ & --- & $+6.3\%$ & $+4.9\%$ \\
$b{=}3$ & --- & $+3.8\%$ & $+4.3\%$ \\
$b{=}4$ & --- & $+4.0\%$ & $+4.8\%$ \\
\bottomrule
\end{tabular}
\end{table}

FajarQuant consistently achieves the lowest key MSE at all bit widths, with 4--6\% improvement over both baselines. The improvement is modest but \textit{consistent}---PCA rotation provides a reliable edge on data with moderate low-rank structure.

% ═══════════════════════════════════════════════════════════
\section{Innovation 2: Fused Quantized Attention}

Standard attention dequantizes all $N$ key vectors before computing dot products, requiring $O(N \cdot d)$ extra memory. We eliminate this by computing directly on quantized indices via codebook lookup:
\begin{equation}
\text{score}_i = q^T \cdot c[k_{\text{idx}}[i]], \quad
\text{output} = \sum_i w_i \cdot c[v_{\text{idx}}[i]]
\end{equation}

\begin{theorem}[Fused Attention Equivalence]
The fused computation produces results identical to standard dequantize-then-attend, with extra memory $O(2^b)$ instead of $O(N \cdot d)$.
\end{theorem}

\begin{table}[t]
\centering
\small
\caption{Memory savings with fused attention ($d{=}128$, $b{=}3$)}
\label{tab:memory}
\begin{tabular}{lccc}
\toprule
Seq.\ Length & FP64 KV & Quantized & Fused Extra \\
\midrule
1K tokens & 2,048 KB & 256 KB & 64 B \\
4K tokens & 8,192 KB & 1,024 KB & 64 B \\
16K tokens & 32,768 KB & 4,096 KB & 64 B \\
\bottomrule
\end{tabular}
\end{table}

% ═══════════════════════════════════════════════════════════
\section{Innovation 3: Hierarchical Multi-Resolution}

We define a bit schedule assigning more bits to recent tokens:

\begin{table}[t]
\centering
\small
\caption{Hierarchical bit budget ($N{=}10{,}000$, $d{=}128$)}
\label{tab:budget}
\setlength{\tabcolsep}{5pt}
\begin{tabular}{lcccc}
\toprule
Tier & Token Age & Bits & Count & Subtotal \\
\midrule
1 & 0--256 & 4 & 256 & 1,024 \\
2 & 257--1024 & 3 & 768 & 2,304 \\
3 & 1025--4096 & 2 & 3,072 & 6,144 \\
4 & 4097+ & 1 & 5,904 & 5,904 \\
\midrule
\textbf{Total} & --- & \textbf{1.54 avg} & 10,000 & \textbf{15,376} \\
Flat 3-bit & --- & 3.0 & 10,000 & 30,000 \\
\textbf{Savings} & & & & \textbf{48.7\%} \\
\bottomrule
\end{tabular}
\end{table}

% ═══════════════════════════════════════════════════════════
\section{Compile-Time Safety}

Fajar Lang's dual-context annotation system provides safety guarantees unavailable in existing frameworks:

\begin{lstlisting}[language=C,caption={@kernel ensures allocation-free quantization}]
@kernel fn quantize(x: [f32; 128],
    codebook: [f32; 8]) -> [u8; 128] {
  // GUARANTEED: no heap, no tensor, no GC
  let mut idx = [0u8; 128]
  for i in 0..128 {
    idx[i] = nearest_centroid(x[i], codebook)
  }
  idx
}
\end{lstlisting}

% ═══════════════════════════════════════════════════════════
\section{Cross-Architecture Perplexity Evaluation}
\label{sec:perplexity}

We evaluate perplexity on WikiText-2 using the R-$\alpha$.1 model
surgery protocol (\S\ref{sec:methodology}) across three models with
fundamentally different KV-head architectures. Five methods are
compared: FajarQuant v1, FajarQuant v2, KIVI, TurboQuant (with
outlier handling), and the FP16 baseline.

\begin{table*}[t]
\centering
\small
\caption{WikiText-2 perplexity (R-$\alpha$.1 protocol, $L{=}2048$,
$n{=}30$ chunks, 61,410 tokens). Lower is better. Bold = best
quantized method per cell. Data source:
\texttt{data/kv\_cache/perplexity\_v2\_\{gemma,mistral,qwen2\}.json}.}
\label{tab:ppl_crossmodel}
\begin{tabular}{ll|ccccc}
\toprule
Model & Bits & FP16 & FQ v1 & \textbf{FQ v2} & KIVI & TQ outlier \\
\midrule
\multirow{3}{*}{Gemma 4 E2B (MQA, 1 KV head)}
 & 2 & 28.13 & 125.19 & 46.20 & 470.50 & \textbf{39.73} \\
 & 3 & 28.13 & 24.26 & 24.55 & \textbf{21.90} & 26.96 \\
 & 4 & 28.13 & 26.75 & 26.98 & 35.17 & 27.33 \\
\midrule
\multirow{3}{*}{Mistral 7B (GQA, 8 KV heads)}
 & 2 & 5.67 & 868.57 & 208.16 & \textbf{23.96} & 167.89 \\
 & 3 & 5.67 & 19.44 & 9.52 & \textbf{5.99} & 9.48 \\
 & 4 & 5.67 & 6.07 & 5.88 & \textbf{5.73} & 5.88 \\
\midrule
\multirow{3}{*}{Qwen2-7B (GQA, 4 KV heads)}
 & 2 & 7.55 & 262.89 & 50.10 & \textbf{46.70} & 165.60 \\
 & 3 & 7.55 & 8.93 & 8.21 & \textbf{8.01} & 8.26 \\
 & 4 & 7.55 & 7.74 & 7.67 & \textbf{7.62} & 7.67 \\
\bottomrule
\end{tabular}
\end{table*}

\subsection{Architecture Dependence Is the Dominant Signal}

The most striking finding is that \textbf{no single method wins
everywhere}. KIVI wins 6 of 9 cells, but catastrophically fails on
Gemma MQA at 2-bit (PPL~470, $16.7\times$ FP16). This is because
KIVI's per-channel symmetric quantization has no redundancy to exploit
when there is only 1~KV head.

FajarQuant v2 is the \textbf{strongest non-KIVI method}, placing 2nd
in 4 of 9 cells. On Qwen2-7B at 2-bit, FQ~v2 achieves PPL~50.10 vs
TQ~outlier's 165.60 (\textbf{69.7\% better}), approaching KIVI's
46.70. The ``right'' quantization method depends on the KV head count:

\begin{itemize}[nosep]
\item \textbf{MQA (1 head):} Rotation-based methods essential. KIVI
      fails catastrophically.
\item \textbf{Wide-GQA ($\geq$8 heads):} KIVI's simplicity wins.
      Rotation adds overhead without benefit.
\item \textbf{Narrow-GQA (4 heads):} In between. FQ~v2 competitive
      with KIVI.
\end{itemize}

\subsection{FQ v2 vs v1: The Upgrade Is Validated}

FQ~v2 improves over v1 in 7 of 9 cells, with massive gains at 2-bit:
63.1\% on Gemma, 76.0\% on Mistral, 80.9\% on Qwen2. The calibrated
PCA + outlier extraction redesign addresses the fundamental weakness
of v1's naive per-tensor quantization.

% ═══════════════════════════════════════════════════════════
\section{Cross-Model Validation}
\label{sec:crossmodel}

To verify that FajarQuant's PCA rotation advantage is not specific to
Gemma~4~E2B's MQA architecture, we extract real KV cache from two
additional models with fundamentally different attention designs and
re-run the 3-way comparison (Table~\ref{tab:crossmodel}).

\begin{table*}[t]
\centering
\small
\caption{Cross-architecture key MSE delta (\% — positive = FajarQuant
better, negative = baseline better). FajarQuant wins keys across all
3 architectures and all 3 bit widths (9 of 9 cells), confirming the
PCA-rotation advantage is architecture-agnostic. Both Mistral 7B and
Qwen2-7B data are from 50-prompt extractions: Mistral $n{=}12{,}800$
samples per cell, Qwen2-7B $n{=}5{,}600$. Gemma~4~E2B numbers are
from Table~\ref{tab:ablation} ($n{=}1{,}500$). The KV-head ratio
column reports attention\_heads : KV\_heads, where $1{:}1$ is full
multi-head attention.}
\label{tab:crossmodel}
\begin{tabular}{lccccccc}
\toprule
\multirow{2}{*}{Model} & \multirow{2}{*}{Arch.} & \multirow{2}{*}{Layers} & \multirow{2}{*}{Heads (a:k)} & \multirow{2}{*}{$\sigma_1/\sigma_2$} & \multicolumn{3}{c}{Key MSE $\Delta$ vs TurboQuant} \\
\cmidrule(lr){6-8}
 & & & & & 2-bit & 3-bit & 4-bit \\
\midrule
Gemma~4~E2B  & MQA      & 35 & $8{:}1$ & 2.67 & $\mathbf{+4.9\%}$ & $\mathbf{+4.3\%}$ & $\mathbf{+4.8\%}$ \\
Mistral 7B   & GQA      & 32 & $32{:}8$ & 2.95 & $\mathbf{+3.6\%}$ & $\mathbf{+4.2\%}$ & $\mathbf{+4.3\%}$ \\
Qwen2-7B     & GQA      & 28 & $28{:}4$ & $\mathbf{7.71}$ & $\mathbf{+4.7\%}$ & $\mathbf{+5.4\%}$ & $\mathbf{+5.4\%}$ \\
\bottomrule
\end{tabular}
\end{table*}

\textbf{Key finding:} FajarQuant wins keys at all 3 bit widths on all
3 architectures, with single-digit-percent improvements that stay
remarkably consistent across model families. The PCA-rotation
advantage is \emph{architecture-agnostic}: it generalizes from MQA
(Gemma's 1 KV head shared across 8 attention heads) through GQA at
both moderate ($4{:}1$, Mistral) and aggressive ($7{:}1$, Qwen2)
sharing ratios. The smallest delta in the table ($+3.6\%$ for Mistral
2-bit vs TurboQuant at $n{=}12{,}800$) is well above any
statistical-noise threshold and confirms the conclusion at scale.

% ═══════════════════════════════════════════════════════════
\subsection{When the PCA Value Gap Closes}
\label{sec:value-gap-closes}

A consistent pattern in the Gemma~4~E2B and Mistral 7B value-MSE
results is that TurboQuant beats FajarQuant by $3$--$11\%$ at every
bit width on values, even though FajarQuant wins keys. This is the
expected design tradeoff for any rotation method optimized on key
distributions: PCA rotation aligns coordinates with the principal
components of the \emph{key} calibration buffer, which need not match
the value distribution.

\textbf{Qwen2-7B exhibits a regime change.} On Qwen2-7B at 2-bit
quantization, the FajarQuant-vs-TurboQuant value-MSE delta closes
entirely to $+0.3\%$ (statistical tie at $n{=}5{,}600$), versus
$-2.8\%$ on Gemma and $-11.1\%$ on Mistral at the same bit width.
This is the first model in our evaluation where FajarQuant's value
performance matches TurboQuant's at the most aggressive quantization
level. The tie was observed at both $n{=}560$ (5-prompt dry run,
$+0.2\%$) and $n{=}5{,}600$ (50-prompt production, $+0.3\%$),
confirming it is not a small-sample artifact.

The closure correlates with Qwen2-7B's exceptional key low-rank
structure: its top-2 singular value ratio is $\sigma_1/\sigma_2 = 7.70$
in Table~\ref{tab:crossmodel}, versus $2.67$ for Gemma and $2.95$ for
Mistral. We hypothesize that when keys are highly low-rank, the PCA
rotation matrix that FajarQuant computes from the key calibration
buffer also produces a useful rotation for the value path, because
both share the same spatial directions in the residual stream of a
transformer. On flatter-distribution models, the key-derived rotation
is less informative for values and TurboQuant's data-agnostic random
rotation wins.

This finding suggests two practical implications:
\begin{enumerate}[nosep]
\item For models with strongly low-rank keys (high
$\sigma_1/\sigma_2$), FajarQuant is competitive on \emph{both} keys
and 2-bit values---the value tradeoff documented elsewhere in this
paper does not apply.
\item Future work could compute a separate value-distribution PCA
rotation cache, decoupling the value-side rotation from the key
calibration. Preliminary experiments are out of scope for this paper,
but the Qwen2-7B result suggests the upper bound is at least
``competitive with TurboQuant on values.''
\end{enumerate}

% ═══════════════════════════════════════════════════════════
\section{KV Cache Structure Analysis}

We analyze the singular value distribution of real KV cache from Gemma~4~E2B (Table~\ref{tab:structure}).

\begin{table}[t]
\centering
\small
\caption{KV cache structure: real Gemma~4 vs synthetic}
\label{tab:structure}
\begin{tabular}{lcc}
\toprule
Metric & Real Gemma~4 & Synthetic \\
\midrule
Key $\sigma_1/\sigma_2$ ratio & $2.73 \pm 0.56$ & $1.11 \pm 0.07$ \\
Key top-3 energy & $66.7\%$ & $99.7\%$ \\
Value $\sigma_1/\sigma_2$ ratio & $1.69 \pm 0.31$ & --- \\
\midrule
Per-layer range & 2.04--4.04 & --- \\
Samples & 150 & 50 \\
\bottomrule
\end{tabular}
\end{table}

Real KV cache keys exhibit $2.45\times$ stronger low-rank structure than synthetic random data. The structure varies significantly across layers: Layer~1 shows the strongest low-rank pattern ($\sigma_1/\sigma_2 = 4.04$), while Layer~9 is most uniform ($\sigma_1/\sigma_2 = 2.04$). This per-layer variation motivates adaptive per-head calibration.

% ═══════════════════════════════════════════════════════════
\section{Ablation Study}

We isolate the contribution of each innovation on real Gemma~4~E2B key data (150 samples, 10 prompts). Table~\ref{tab:ablation} shows incremental gains as innovations are added.

% Full-width table to fit all 5 columns without truncation
\begin{table*}[t]
\centering
\small
\caption{Ablation: incremental innovation contributions (Gemma~4~E2B, $N{=}10$K)}
\label{tab:ablation}
\begin{tabular}{lcccc}
\toprule
Configuration & Key MSE ($b{=}2$) & MSE $\Delta$ vs TurboQuant & Dequant Memory & Bit Savings \\
\midrule
No rotation (KIVI-like) & $9.31 \times 10^{-4}$ & $+1.5\%$ & $O(N \cdot d)$ & 0\% \\
Random rotation (TurboQuant) & $9.18 \times 10^{-4}$ & --- & $O(N \cdot d)$ & 0\% \\
\midrule
$+$ PCA rotation & $\mathbf{8.72 \times 10^{-4}}$ & $-4.9\%$ & $O(N \cdot d)$ & 0\% \\
$+$ Fused attention & $\mathbf{8.72 \times 10^{-4}}$ & $-4.9\%$ & $O(2^b)$ & 0\% \\
$+$ Hierarchical (Full FajarQuant) & $\mathbf{8.72 \times 10^{-4}}$ & $-4.9\%$ & $O(2^b)$ & $48.7\%$ \\
\bottomrule
\end{tabular}
\end{table*}

\textbf{Key findings:} (1)~PCA rotation provides 4--5\% MSE improvement across all bit widths (4.9\% at 2-bit, 4.3\% at 3-bit, 4.8\% at 4-bit), confirming that exploiting low-rank structure in KV cache is consistently beneficial. (2)~Fused attention preserves identical MSE while reducing dequantization memory from $O(N \!\cdot\! d)$ to $O(2^b)$---a $524{,}288\times$ reduction at $N{=}16$K. (3)~Hierarchical allocation saves 48.7\% of bit budget at $N{=}10$K by assigning 4~bits to the 256 most recent tokens and 1~bit to tokens beyond position 4096. The savings increase with context length (55.7\% at $N{=}16$K).

Notably, ``no rotation'' (equivalent to KIVI per-channel) is 1.5\% \textit{worse} than random rotation at 2-bit, confirming that rotation-based approaches---whether random or adaptive---are fundamentally superior for KV cache quantization.

% ═══════════════════════════════════════════════════════════
\section{End-to-End Results}

\begin{table}[t]
\centering
\small
\caption{End-to-end summary}
\label{tab:e2e}
\setlength{\tabcolsep}{4pt}
\begin{tabular}{llr}
\toprule
Contribution & Metric & Result \\
\midrule
Cross-arch analysis & Models tested & 3 (MQA+GQA) \\
FQ v2 vs v1 (2-bit) & PPL improvement & 63--81\% \\
FQ v2 vs TQ (Qwen2) & PPL at 2-bit & $\mathbf{-69.7\%}$ \\
Fused ($N{=}16$K) & Mem.\ saved & 4\,MB$\to$64\,B \\
Hierarchical & Bit savings & $\mathbf{48.7\%}$ \\
Native impl & vs Python & $\mathbf{5.0\times}$ faster \\
Safety & Overhead & Zero \\
\bottomrule
\end{tabular}
\end{table}

\subsection{Memory Budget}

\begin{table}[t]
\centering
\small
\caption{End-to-end memory ($L{=}32$, $H{=}32$, $d{=}128$)}
\label{tab:memory_e2e}
\setlength{\tabcolsep}{4pt}
\begin{tabular}{lcccc}
\toprule
$N$ & FP16 & Flat 3-bit & Hier. & Ratio \\
\midrule
1K & 512 MB & 64 MB & 35 MB & 14.6$\times$ \\
4K & 2 GB & 256 MB & 140 MB & 14.6$\times$ \\
16K & 8 GB & 1 GB & 547 MB & 14.6$\times$ \\
64K & 32 GB & 4 GB & 2.19 GB & 14.6$\times$ \\
\bottomrule
\end{tabular}
\end{table}

% ═══════════════════════════════════════════════════════════
\section{Related Work}

\textbf{TurboQuant} \citep{zandieh2025turboquant}: Our baseline. Achieves $2.7\times$-optimal MSE via random rotation + Lloyd-Max.

\textbf{KIVI} \citep{liu2024kivi}: Per-channel key / per-token value quantization. Strong for key MSE but we show rotation-based methods are superior overall, especially at low bit widths.

\textbf{AQLM} \citep{egiazarian2024aqlm}: Additive quantization for model weights (not KV cache).

\textbf{QuIP\#} \citep{tseng2024quip}: Incoherence processing + LDLQ for weight quantization.

\textbf{FlexGen} \citep{sheng2023flexgen}: Memory offloading system. Complementary to quantization.

No prior work combines adaptive rotation, fused attention, hierarchical allocation, and compile-time safety.

% ═══════════════════════════════════════════════════════════
\section{Performance Evaluation}
\label{sec:perf}

This section reports wall-clock latency, throughput, and peak memory
for FajarQuant vs the KIVI \citep{liu2024kivi} and TurboQuant
\citep{zandieh2025turboquant} baselines on Lenovo Legion Pro
(Intel i9-14900HX, RTX 4090 Laptop). All numbers are measured under
the protocol locked in
\href{https://github.com/fajarkraton/fajarquant/blob/main/bench/METHODOLOGY.md}{\texttt{bench/METHODOLOGY.md}}
on the hardware captured in
\href{https://github.com/fajarkraton/fajarquant/blob/main/bench/hardware\_snapshot.txt}{\texttt{bench/hardware\_snapshot.txt}}.
The criterion benchmark harness uses 100 samples per measurement,
3-second warmup, 10-second measurement window, and reports the
median plus the 95\% confidence interval (not the mean) for outlier
robustness.

The reproducibility note (\S\ref{sec:reprod}) explains how to re-run
every measurement in this section from a clean checkout.

\subsection{Per-Layer Quantization Latency}
\label{subsec:perf-latency}

We measure the time to quantize one $(seq\_len{=}128, d{=}128)$
layer of a synthetic Mistral-7B-shaped tensor at each bit width.
Synthetic data is appropriate for latency: the per-element
arithmetic and memory traffic do not depend on the input
distribution. Numerical accuracy uses real KV cache (\S\ref{sec:perplexity}).

\begin{table}[h]
\centering
\caption{Per-layer quantization latency (median $\pm$ 95\% CI half-width
over 100 samples). Lower is better. Data source:
\texttt{bench/results/quant\_latency.csv} (V26 Phase C2.1, generated by
\texttt{cargo bench --bench quant\_latency}).}
\label{tab:latency}
\begin{tabular}{l|c|c|c}
\hline
Bit width & KIVI & TurboQuant & \textbf{FajarQuant} \\
\hline
2-bit & TBD & TBD & TBD \\
3-bit & TBD & TBD & TBD \\
4-bit & TBD & TBD & TBD \\
\hline
\end{tabular}
\end{table}

\textit{Status:} the bench harness is committed
(\texttt{benches/quant\_latency.rs}, V26 commit \texttt{5741753}) and
compiles cleanly. Actual measurements pending — they require
\texttt{bash bench/setup\_perf.sh} to be run first to flip CPU
governor $\to$ performance, disable Turbo, take HT siblings offline,
and pin to even cores. The C2.0.3 noise floor capture
(\texttt{bench/results/noise\_floor.json}) shows that without
\texttt{setup\_perf.sh}, sub-nanosecond benchmarks have a 64\% cross-run
drift on this machine.

\subsection{Throughput: Tokens per Second vs FP16}
\label{subsec:perf-throughput}

We measure end-to-end inference throughput at batch sizes 1, 4, and
16 with the KV cache quantized vs the FP16 baseline. Higher is better.

\begin{table}[h]
\centering
\caption{Inference throughput (tokens/sec, median over 100 samples).
Data source: \texttt{bench/results/throughput.csv}.}
\label{tab:throughput}
\begin{tabular}{l|c|c|c|c}
\hline
Configuration & FP16 baseline & KIVI & TurboQuant & \textbf{FajarQuant} \\
\hline
Mistral-7B, batch 1, 2-bit  & TBD & TBD & TBD & TBD \\
Mistral-7B, batch 4, 2-bit  & TBD & TBD & TBD & TBD \\
Mistral-7B, batch 16, 2-bit & TBD & TBD & TBD & TBD \\
\hline
\end{tabular}
\end{table}

\textit{Status:} pending V26 Phase C2.2 substantive work.

\subsection{Peak Memory at 16K Context}
\label{subsec:perf-memory}

We measure peak resident set size (RSS) for a single 16K-context
Mistral-7B inference run with KV cache quantized at each bit width.
This validates the memory savings claimed in \S\ref{sec:fused}
($O(N \cdot d) \to O(2^b)$ extra dequant memory).

\begin{table}[h]
\centering
\caption{Peak RSS at 16K context (megabytes; lower is better).
Data source: \texttt{bench/results/memory\_profile.csv}.}
\label{tab:memory}
\begin{tabular}{l|c|c|c|c}
\hline
Algorithm & 2-bit & 3-bit & 4-bit & FP16 baseline \\
\hline
KIVI       & TBD & TBD & TBD & --- \\
TurboQuant & TBD & TBD & TBD & --- \\
\textbf{FajarQuant} & \textbf{TBD} & \textbf{TBD} & \textbf{TBD} & --- \\
FP16 reference & --- & --- & --- & TBD \\
\hline
\end{tabular}
\end{table}

\textit{Status:} pending V26 Phase C2.3 substantive work. Capture
will use \texttt{valgrind --tool=massif} for ground truth plus a
criterion group that records \texttt{peak\_rss\_kb} via
\texttt{getrusage(RUSAGE\_SELF)}.

\subsection{Wall-Clock vs TurboQuant on Identical Hardware}
\label{subsec:perf-wallclock}

A head-to-head end-to-end run of FajarQuant vs TurboQuant on the same
Mistral-7B input distribution. Reports the median speedup factor with
95\% CI.

\textit{Status:} pending V26 Phase C2.4 substantive work; output
artifact \texttt{paper/figures/wallclock\_bar.pdf}.

\subsection{Methodology Notes}

The benchmark protocol locked in
\texttt{bench/METHODOLOGY.md} \S3 mandates that all wall-clock
measurements be taken on a machine where:

\begin{enumerate}[nosep]
\item CPU governor is pinned to \texttt{performance} (default Ubuntu
24.04 is \texttt{powersave}, which adds 100--300\% timing variance);
\item Intel Turbo Boost is disabled (frequency drift is the dominant
noise source on i9-14900HX);
\item Intel Hyperthreading siblings on each test core are taken
offline (HT siblings share L1+L2 cache);
\item Single-threaded evaluation (\texttt{cargo bench
--test-threads=1}); and
\item Page caches dropped + swap disabled before each run.
\end{enumerate}

The companion script \texttt{bench/setup\_perf.sh} (\texttt{apply}
mode) mechanizes all five steps in a single \texttt{sudo}-prompting
invocation. Without it, the noise floor on this machine is
\textbf{0.720~ns} for sub-nanosecond operations
(\texttt{bench/results/noise\_floor.json}, captured 2026-04-11),
which corrupts measurements smaller than $\sim$1.4~ns. FajarQuant's
per-layer quantize is in the microsecond regime, so the noise floor
matters only when comparing two algorithms whose deltas are
sub-nanosecond---a situation that does not arise in any of the
tables in this section.

% ═══════════════════════════════════════════════════════════
\section{Conclusion}

We present the first systematic cross-architecture evaluation of KV
cache quantization methods, testing five methods across three
fundamentally different attention designs (MQA, wide-GQA, narrow-GQA).
Our primary finding is that \textbf{no single method dominates}: KIVI
wins 6/9 cells but catastrophically fails on MQA at 2-bit, while
FajarQuant~v2 is the strongest non-KIVI method on MQA and narrow-GQA.
The ``right'' method depends on the KV head count.

FajarQuant~v2---combining calibrated PCA rotation with outlier
extraction---improves 63--81\% over v1 at 2-bit, validating the
algorithm redesign. Combined with fused Hadamard+quantize (1.6$\times$
speedup), hierarchical bit allocation (48.7\% savings), and a native
Fajar Lang implementation (5$\times$ faster than Python), the system
enables practical embedded deployment with compile-time safety
guarantees.

\subsection*{Future Work}
\begin{enumerate}[nosep]
\item Adaptive method selection based on detected KV-head architecture.
\item Separate value-distribution PCA rotation (Qwen2-7B results suggest this closes the value gap).
\item Online PCA update during autoregressive generation.
\item Extension from KV cache to model weight quantization.
\item Evaluation on larger models (Gemma~4~31B, Llama~3).
\end{enumerate}

% ═══════════════════════════════════════════════════════════
% Force page break before appendix to prevent overlap
\clearpage
\appendix

\section{Reproducibility}

\begin{lstlisting}[language=bash,caption={Reproducing all results}]
git clone https://github.com/fajarkraton/fajar-lang
cd fajar-lang && cargo build --release

# Rust tests (23 FajarQuant tests)
cargo test --lib -- runtime::ml::fajarquant

# KV cache extraction (requires GPU + PyTorch)
python scripts/extract_kv_cache.py \
  --model google/gemma-4-E2B --num-prompts 50

# 3-way MSE comparison
python scripts/run_comparison.py \
  --data data/kv_cache/ --num-prompts 10

# Perplexity evaluation
python scripts/eval_perplexity.py \
  --model google/gemma-4-E2B --bits 2,3,4

# Ablation study
python scripts/run_ablation.py \
  --data data/kv_cache/ --num-prompts 10
\end{lstlisting}

\textbf{Environment:} NVIDIA RTX 4090 Laptop GPU (16~GB), Intel i9-14900HX, 32~GB RAM, Ubuntu 24.04, Rust~1.87, PyTorch~2.11+cu130, Transformers~5.5.

\textbf{Model:} Google Gemma~4~E2B (5B parameters, 35 layers, 1~KV head, $d{=}256$, Apache~2.0 license).

\textbf{Dataset:} WikiText-2-raw-v1 \citep{merity2017pointer} test split (1.29M characters).

\section{Proof of Theorem~3}

\begin{proof}
The MSE of per-coordinate quantization after rotation $R$ is $\text{MSE}(R) = \sum_i D_i(\sigma_i^2, b)$ where $\sigma_i^2$ is the variance of coordinate $i$ after rotation, and $D_i$ is the per-coordinate distortion.

For PCA rotation $R_{\text{pca}}$, coordinates are principal components with ordered variances $\lambda_1 \geq \cdots \geq \lambda_d$. Per-coordinate quantization allocates independent scale/zero-point per dimension, adapting to each component's dynamic range.

For random rotation $R_{\text{rand}}$, coordinate variances are approximately equalized at $\text{tr}(C)/d$.

With per-coordinate quantization, distortion $D_i \propto \text{range}_i^2 / 4^b$. PCA minimizes total range by concentrating variance: low-variance components have small range and low distortion. The improvement is:
\begin{equation}
\frac{\text{MSE}(R_{\text{pca}})}{\text{MSE}(R_{\text{rand}})} = \frac{\sum_i \lambda_i^2}{\left(\sum_i \lambda_i\right)^2 / d} \leq 1
\end{equation}
with equality iff all eigenvalues are identical (no structure to exploit).
\end{proof}

% ═══════════════════════════════════════════════════════════
\balance
\bibliographystyle{plainnat}
\bibliography{references}

\end{document}
